\section{Introduction}
\label{sec:intro}

Language-model agents are often trained to say what they expect before they act. ReAct-style thoughts, belief states and plans all put a prediction in writing \citep{yao2023react,wang2026evu,zhang2025earlyexperience}. Monitors read these predictions to judge why an agent acts \citep{korbak2025monitorability}, and editing them is a cheap way to steer an agent. Both practices assume the written prediction takes part in the decision.

Work on chain-of-thought faithfulness shows that whether this assumption holds depends on the test. Deleting reasoning often leaves answers unchanged, while editing it can change them \citep{lanham2023faithfulness}. In fine-tuned agents, replacing the whole reasoning trace moves fewer actions as training goes on \citep{liu2026cotgains}.

For knowledge that the model has learned, the question is harder. In ordinary training data the written prediction and the knowledge in the weights rarely conflict, so the data do not show which of the two the action follows when they do.

We study the question in a setting built so that the two can be told apart. \alfcf{} is a counterfactual version of ALFWorld in which each object class is kept in a fixed, arbitrary receptacle class, chosen so that neither the original data nor the base model's prior points to it. The location can therefore only be learned from fine-tuning.

At the first search step of an episode, the agent has not yet seen the target, so its one-line prediction, which we call the \emph{note}, is the only location information in its context (Figure~\ref{fig:overview}). We fine-tune Qwen3-4B on six versions of the training data. They differ in one respect each: whether predictions are always followed by the action, sometimes contradicted by it, or followed by an action that ignores them.

At the first search step we then either delete the agent's prediction or replace it with the prediction the agent writes for another object, and record where the agent goes. Each arm is compared with what a perfect imitator of its training data would do, so effects that follow from the data construction are not read as properties of the model.

% Conceptual overview (no bookkeeping numbers). Included from the introduction.
\begin{figure*}[t]
\centering
\begin{tikzpicture}[
  font=\small,
  box/.style={draw=black!55, rounded corners=3pt, align=left, inner sep=4pt, fill=white, minimum height=1.55cm},
  ctx/.style={box, fill=black!4},
  note/.style={box, fill=blue!7, draw=blue!55!black},
  act/.style={box, fill=black!4},
  wts/.style={draw=black!45, rounded corners=3pt, align=center, inner sep=3pt, fill=black!4, minimum height=0.6cm},
  ivn/.style={box, fill=orange!9, draw=orange!70!black, minimum height=1.3cm},
  read/.style={box, fill=white, draw=black!35, minimum height=1.3cm},
  data/.style={box, fill=green!6, draw=green!45!black},
  arr/.style={-{Latex[length=2.2mm]}, thick, black!65},
  darr/.style={-{Latex[length=2.2mm]}, thick, dashed, orange!75!black},
]
\node[data, text width=3.0cm] (data) {\textbf{Training data}\\[1pt] predictions always followed, sometimes contradicted by the action, or ignored by it};
\node[ctx, right=0.9cm of data, text width=2.6cm] (ctx) {\textbf{Context}\\[1pt] the room and the task; the target is not yet in view};
\node[note, right=0.4cm of ctx, text width=3.3cm] (note) {\textbf{Written prediction}\\[1pt] \texttt{The mug is most likely in a toaster, then \dots}};
\node[act, right=0.4cm of note, text width=1.9cm] (act) {\textbf{Next action}\\[1pt] \texttt{go to ?}};
\node[wts, above=0.55cm of note.north west, anchor=south west, text width=6.3cm] (wts) {\textbf{Model weights}: location knowledge from fine-tuning};
\draw[arr, green!45!black] (data.north) |- node[pos=0.72, above, font=\scriptsize, green!35!black] {fine-tune} (wts.west);
\draw[arr] (ctx) -- (note);
\draw[arr] (note) -- (act);
\draw[arr, black!40] (wts.south -| note.north) -- (note.north);
\draw[arr, black!40] (wts.south -| act.north) -- (act.north);
\node[ivn, below=0.6cm of note, text width=3.3cm] (ivn) {\textbf{Interventions}\\[1pt] \emph{flip}: the agent's own prediction for another object\\ \emph{remove}: no prediction};
\draw[darr] (ivn) -- (note);
\node[read, below=0.6cm of act.south west, anchor=north west, text width=3.7cm] (read) {\textbf{Reading}\\[1pt] the action follows the flip: the prediction steers the step\\ the action survives removal: the prediction is not needed};
\draw[arr, black!40] (act.south) -- (act.south |- read.north);
\end{tikzpicture}
\caption{The question and the design. Before each search step the agent writes where it expects the target. Fine-tuning puts the location knowledge into the weights, from which it can reach the action directly or through the written prediction. We fine-tune on data that differ only in how predictions and actions are coupled. At the first search step, where the prediction is the only location information in context, we flip or remove the prediction and record where the agent goes. Each training arm is also compared with what a perfect imitator of its data would do.}
\label{fig:overview}
\end{figure*}


When every prediction in the training data is followed, the prediction steers the agent without being needed. Deleting the prediction lowers the probability of the predicted move by 0.015. Replacing it redirects 80\% of first moves across eight seeds, and a randomly ordered ranking, a form never seen in training, is followed even more readily (0.90). In held-out rooms, whether the weights hold a convention for the object makes no detectable difference.

Over a whole episode, writing another object's prediction at every step raises the number of receptacles visited from about two to about nine. The control is strongest at the first step and much weaker later, where it depends on what the search has already shown. Removing the agent's earlier predictions from its history partly restores the control.

When a quarter or half of the training predictions are contradicted by the teacher's action, the agents keep writing predictions but follow a replaced one in at most about 1\% of first steps, within 0.012 of what a perfect imitator of such data does. When half of the actions ignore correct predictions instead, redirection drops to 27\%, between what an imitator would show if it credited its correct actions to the weights (0) and to the predictions (0.50).

For monitoring, this means that when the training data also teach the agent to act without a prediction, as ours do at 30\% of steps, a deletion test calls the predictions decorative, while an edit test shows that they steer the first decision. For training, contradicting a quarter of the predictions was enough to remove that steering.
